% =====================================================================
% Premessa
% =====================================================================
% copre: docs/31-glossario.md
% copre: SOURCES.md#come-si-usa
% copre: SOURCES.md#due-file-due-scopi
% verificato-al-commit: 3bc311b
% ---------------------------------------------------------------------

\section*{Origine del lavoro}

Questo lavoro nasce da una cartuccia guasta e da una curiosità rimasta senza risposta per
vent'anni.

La cartuccia è una copia originale di Pokemon Smeraldo, giocata da bambino e conservata,
la cui tasca degli strumenti a un certo punto ha smesso di avere senso: da uno slot in
avanti gli oggetti ordinari risultano sostituiti da Poke Ball. Non è un guasto che si
risolva ricominciando la partita, perché quel salvataggio contiene uno stato che non si
riproduce. Da qui la prima domanda, concreta e non teorica: se sia possibile estrarre un
salvataggio di vent'anni, determinare che cosa lo abbia alterato e ripararlo senza
compromettere il resto.

La seconda origine è più antica della prima. I giochi Pokemon delle prime tre generazioni
comunicano fra loro attraverso un cavo, e ciò che quel cavo trasporta sono creature. Un
Pokemon che passa da una console a un'altra è, sotto la superficie, una struttura di
ottanta byte che attraversa un canale seriale sincrono, viene cifrata, verificata da un
checksum e ricomposta all'altro capo. La formazione in ingegneria delle telecomunicazioni
fornisce il vocabolario per descrivere quel meccanismo; l'interesse per il dominio fornisce
la pazienza necessaria a leggerne i disassemblati. Questo documento nasce dalla
combinazione delle due cose.

Da queste premesse sono derivati cinque lavori paralleli, che costituiscono i casi di
studio della trattazione. Il primo è la riparazione dell'inventario della cartuccia
citata, per estrazione fisica del salvataggio. Il secondo è la costruzione di un ponte che
trasferisca esemplari dalle prime due generazioni alla terza su hardware originale,
funzione che il produttore non ha mai reso disponibile e che va quindi realizzata. Il terzo
è la modifica di un Nintendo 3DS per estrarre i dati delle cartucce possedute. Il quarto è
la comunicazione fra un calcolatore e una Nintendo Switch attraverso il protocollo
wireless locale proprietario, per eseguire uno scambio con un gioco in esecuzione sulla
console. Il quinto è lo studio dell'automazione di un gioco mediante pilotaggio del
controller e analisi del segnale video.

Nessuno dei cinque è nato come esercizio: ciascuno è nato da un obiettivo pratico
irraggiungibile senza comprendere il meccanismo sottostante. È questa la ragione per cui
ogni capitolo teorico termina in un vincolo operativo e ogni vincolo operativo risale a un
principio.

\section*{Oggetto della trattazione}

Il testo che segue non è il resoconto dei cinque lavori, ma l'esposizione sistematica di
ciò che è stato necessario apprendere per condurli. L'oggetto è la rappresentazione, la
verifica di integrità, la trasmissione e la conversione di un dato strutturato fra sistemi
digitali non progettati per interoperare. I dati dei giochi costituiscono il carico utile,
le console i nodi.

La scelta del dominio è motivata tecnicamente. Quei giochi sono sistemi reali,
integralmente documentati da decompilazioni fedeli che ricompilano in immagini identiche
agli originali, abbastanza contenuti da poter essere compresi per intero e abbastanza
antichi da non presentare alcuna protezione che offuschi i meccanismi. Costituiscono
pertanto un laboratorio più accessibile della maggior parte dei sistemi contemporanei.

L'esposizione muove dai fondamenti elementari. La prima parte non presuppone la conoscenza
della rappresentazione binaria, dell'aritmetica in complemento, della distinzione fra
indirizzo e valore o della nozione di esecuzione. Il resto del documento vi si appoggia
esplicitamente, e in più punti la spiegazione di un fenomeno complesso si riduce al
richiamo di un fatto elementare stabilito in quella sede.

Dove occorre un concetto più avanzato, esso viene introdotto nel punto in cui serve. Quando
tale concetto possiede una denominazione consolidata nella teoria delle telecomunicazioni,
quella denominazione viene adottata: riconoscere che un problema apparentemente
idiosincratico è un caso particolare di un problema noto costituisce metà della sua
soluzione. Il checksum è trattato come codice di rilevazione d'errore; la conversione delle
statistiche fra due generazioni come quantizzazione con saturazione; lo stallo fra due
protocolli asimmetrici come problema di controllo di flusso; il collegamento fra due
console come canale sincrono privo di controllo d'errore.

\section*{Metodo e gerarchia delle fonti}

Ogni affermazione tecnica riporta la fonte da cui deriva, e le fonti sono ordinate secondo
una gerarchia di affidabilità decrescente che è normativa: quando due fonti si
contraddicono, prevale quella di livello inferiore.

Il \term{livello 1} comprende il codice dei giochi e la documentazione dell'hardware
ricostruita e verificata dalla comunità. Nel caso dei giochi si tratta di decompilazioni
che ricompilano in immagini binariamente identiche alle originali, il che le rende
definizione del comportamento e non descrizione: \file{pret/pokered} per la prima generazione,
\file{pret/pokecrystal} per la seconda, \file{pret/pokeemerald} e \file{pret/pokefirered} per la terza.
Per l'hardware la funzione è svolta da GBATEK per il Game Boy Advance e dalle Pan Docs per
il Game Boy. Questo livello costituisce la verità operativa.

Il \term{livello 2} comprende la documentazione di dominio, cioè wiki e opere di
riferimento. È accurata nell'impianto generale e non infallibile nel dettaglio, il che la
rende adatta a orientarsi e inadatta a fondarvi una riga di codice senza verifica.

Il \term{livello 3} comprende le implementazioni di riferimento, cioè codice di terzi
funzionante sul campo. Tale codice incorpora conoscenza validata dall'uso, e insieme scelte
arbitrarie del suo autore che non vanno confuse con specifiche del formato: dimostra che
una cosa è possibile, non che quello sia il modo corretto di farla.

Il \term{livello 4} comprende articoli, blog, diari di sviluppo e materiale video. È
prezioso per la ricostruzione delle motivazioni progettuali e inaffidabile per i valori
numerici.

Il \term{livello 5} comprende forum e comunità. È spesso l'unica sede in cui una
determinata informazione esista, e va trattato come testimonianza: non è citabile finché
non si verifica su un livello superiore.

\subsection*{La casistica che ha prodotto la gerarchia}

L'ordinamento non è un formalismo metodologico ma il risultato di errori concreti,
riscontrati durante la redazione della referenza dei formati di questo progetto e qui
elencati perché costituiscono la sua giustificazione empirica.

Il livello 2 collocava le cifre della tabella caratteri di prima generazione all'offset
\hx{F0} in luogo di \hx{F6}, e le maiuscole di quella di terza generazione a \hx{C1} in
luogo di \hx{BB}: entrambi gli errori producono testo plausibile e falso, cioè il genere di
difetto che non si manifesta come tale. Descriveva il calcolo del checksum di terza
generazione come somma di byte in luogo di somma di parole da sedici bit, errore
sufficiente a distruggere un esemplare in modo definitivo. Riportava due valori in
conflitto per la dimensione del blocco di scambio, 415 e 424, mentre la somma delle
costanti nel disassemblato stabilisce 424 byte trasmessi sul canale di cui 418 di dati.

Il livello 3 si è rivelato in disaccordo con la propria documentazione: la libreria di
conversione della comunità documenta quattro metodi nel proprio manuale e ne implementa uno,
e la funzione destinata alla conversione dei punti di allenamento contiene un ciclo che
azzera tutti e sei i campi.

In ciascuno di questi casi ha prevalso il livello 1.

\subsection*{Trattamento del non verificato}

Un'affermazione non verificata non viene esposta come verificata. Le proposizioni che
poggiano su inferenza, su testimonianza non confermata o su lettura non ancora ricontrollata
sul sorgente sono racchiuse in un riquadro che le dichiara tali. Il lettore che incontra un
riquadro sa di trovarsi davanti a un debito e non a un risultato.

Questa disciplina ha un'origine documentata. Quattro affermazioni sulle quali si stava per
fondare del codice si sono rivelate errate, ed è emerso soltanto clonando i disassemblati e
cercando nel sorgente. Il clone superficiale è pertanto il primo passo di ogni verifica e
non l'ultimo.

\section*{Rapporto con gli altri documenti del progetto}

Il presente documento è la forma organizzata e integrale del materiale tecnico prodotto nel
corso del lavoro. Le note sotto \file{docs/} costituiscono il percorso di studio, che tratta
gli stessi argomenti per un lettore già in possesso del contesto e in forma navigabile come
grafo. La referenza dei formati documenta le strutture byte per byte e serve alla scrittura
del codice. Le schede sotto \file{.claude/} registrano lo stato di avanzamento del lavoro,
non conoscenza, e non hanno corrispondenza in questa sede.

Tre testi che espongono la stessa materia in registri diversi divergono, in assenza di un
vincolo che lo impedisca. Il vincolo adottato è un controllo che fallisce: ogni capitolo
dichiara quali documenti e quali loro sezioni riorganizza, e a quale revisione del
repository la dichiarazione è stata verificata; lo strumento \file{tools/check-thesis-coverage.py}
segnala le sezioni che nessun capitolo reclama, i capitoli il cui documento di riferimento
è stato modificato dopo l'ultima verifica, e i riferimenti bibliografici privi di voce.

La bibliografia non è redatta a mano: è generata dalla medesima tabella da cui derivano le
note sulle fonti, cosicché le due non possano riportare informazioni discordanti.

\section*{Convenzioni}

I valori esadecimali sono nella forma \hx{2A}, i binari nella forma \bin{00101010}. I
percorsi dei file del progetto sono in \file{questo stile}. Un termine tecnico alla prima
occorrenza è in \term{corsivo}, e i termini raccolti nel glossario del progetto conservano
qui la medesima definizione. Gli estratti di codice riportano in didascalia il file e la
fonte di provenienza, poiché un estratto privo di provenienza non è verificabile.
